\section{\textbf{\texttt{self.env}} Explained}

\begin{frame}[fragile]{What Is \texttt{self.env}?}
	\begin{itemize}
		\item \textbf{\texttt{self.env}} is the \textbf{Environment} object attached to a recordset.
		\item It is the gateway to models, users, companies, context, and the database cursor.
\begin{block}{Basic Access}
\begin{lstlisting}
class Student(models.Model):
	_name = 'student.student'

	def action_demo(self):
		print(self.env)
		students = self.env['student.student'].search([])
\end{lstlisting}
\end{block}
		\item Inside model methods, you almost always start from \texttt{self.env}.
	\end{itemize}
\end{frame}

\begin{frame}[fragile]{Environment vs Recordset}
	\begin{itemize}
		\item A \textbf{recordset} represents one or more records of a model.
		\item An \textbf{environment} is the execution context around those records.
\begin{lstlisting}
student = self.env['student.student'].browse(15)
print(student)      # student.student(15,)
print(student.env)  # Environment object
\end{lstlisting}
		\item Different recordsets may share the same environment.
		\item Changing the environment (for example with \texttt{sudo()}) returns a \textbf{new} recordset / env pair.
	\end{itemize}
\end{frame}

\begin{frame}[fragile]{Where Does \texttt{self.env} Come From?}
	\begin{itemize}
		\item Odoo builds the environment when a request / shell / cron job starts.
		\item Then every recordset created in that flow receives that environment.
\begin{lstlisting}
# In Odoo shell
env['student.student'].search([])

# In a model method
self.env['student.student'].search([])
\end{lstlisting}
		\item In both cases, \texttt{env} is the same kind of object.
		\item In model methods, it is available as \textbf{\texttt{self.env}}.
	\end{itemize}
\end{frame}
